\chapter{ASIC Synthesis}
\label{ch:syn}

The ASIC flow synthesizes one unit with Cadence Genus. It is in \file{lib/syn},
and each unit adds a Makefile and the timing constraints
\file{<unit>/hw/syn/umc130/<unit>.sdc}. The supported node is UMC 130~nm, with
the library setup in \file{lib/syn/umc130/genus.setup.tcl}. The commands and
variables are in Section~\ref{sec:make_syn}.

The script \file{lib/syn/genus/build.tcl} runs these steps in batch mode:

\begin{enumerate}
\item Load the technology library from the directory in \sig{UMC130\_DIR}. The
  flow stops with an error if the variable is not set.
\item Read the sources with the macro \sig{ASIC} defined, and elaborate the unit.
\item Read the constraints and check the design.
\item Run generic synthesis, mapping and optimization.
\item Insert tie cells and remove unloaded and undriven logic.
\item Write the reports and the netlist.
\end{enumerate}

Table~\ref{tab:syn_out} lists the outputs, in \file{<unit>/hw/syn/results}. A
remote run copies the logs, the reports and the netlist back to
\file{<unit>/hw/syn}.

\begin{twotable}{ASIC Synthesis Output Files}{tab:syn_out}{File}{Content}{5.2cm}
\file{<unit>\_synth.v} & Mapped netlist. \\ \hline
\file{<unit>\_check\_design.rpt} & Design check. \\ \hline
\file{<unit>\_lint\_gen\_timing.rpt} & Timing lint after generic synthesis. \\ \hline
\file{<unit>\_area.rpt}, \file{<unit>\_gates.rpt} & Area and gate reports. \\ \hline
\file{<unit>\_clk.rpt}, \file{<unit>\_timing.rpt} & Clock and timing reports. \\ \hline
\file{<unit>\_power.rpt}, \file{<unit>\_qor.rpt} & Power and quality-of-results reports. \\
\end{twotable}

\section{Results}
\label{sec:syn_results}

Table~\ref{tab:syn_results} gives the Genus 19.1 synthesis results in UMC
130~nm, with the Faraday \sig{fsc0l\_d} 2009Q2v3.0 low-leakage library at the
slow corner (1.08~V, 125~\textdegree C), an 8.2~ns clock (122~MHz) and the default
parameters (binary32). The critical period is the clock period minus the worst
slack. The power is estimated with the default switching activity, without
simulation vectors.

\begin{table}[ht]
\centering
\small
\begin{tabular}{|l|r|r|r|r|r|}
\hline
\rowcolor{ug-head}\bfseries Unit & \bfseries Cells & \bfseries Area ($\mu$m$^2$) & \bfseries Slack (ns) & \bfseries Period (ns) & \bfseries Power (mW) \\ \hline\hline
\sig{fadd} & 1,239 & 15,843 & 0.003 & 8.20 & 0.8 \\ \hline
\sig{fmul} & 2,931 & 37,324 & 0.001 & 8.20 & 1.5 \\ \hline
\sig{fdiv} & 1,382 & 15,265 & 0.001 & 8.20 & 0.7 \\ \hline
\sig{fmac} & 7,703 & 84,739 & 0.000 & 8.20 & 8.5 \\ \hline
\sig{fpu} & 5,749 & 67,341 & 0.000 & 8.20 & 2.5 \\ \hline
\end{tabular}
\caption{ASIC Synthesis Results at 122~MHz}
\label{tab:syn_results}
\end{table}

All units meet the 8.2~ns clock. The critical path of \sig{fmul} and \sig{fpu} is
the mantissa multiplication, which takes one pipeline stage. In \sig{fmac} it is the
normalization of the accumulator. In \sig{fadd} it is the rounding in \sig{pack},
and in \sig{fdiv} the exponent of the quotient, computed when the operands are
captured.
